\documentclass[11pt]{article}
\usepackage[margin=1in]{geometry}
\usepackage[T1]{fontenc}
\usepackage[utf8]{inputenc}
\usepackage[french]{babel}
\usepackage{amsmath,amssymb,amsthm,mathtools,bm}
\usepackage{enumitem}
\usepackage{booktabs}
\usepackage{xcolor}
\usepackage[hidelinks]{hyperref}
\usepackage{microtype}

\theoremstyle{plain}
\newtheorem{theoreme}{Théorème}[section]
\newtheorem{proposition}[theoreme]{Proposition}
\newtheorem{lemme}[theoreme]{Lemme}
\newtheorem{corollaire}[theoreme]{Corollaire}
\theoremstyle{definition}
\newtheorem{definition}[theoreme]{Définition}
\newtheorem{exemple}[theoreme]{Exemple}
\theoremstyle{remark}
\newtheorem{remarque}[theoreme]{Remarque}

\DeclareMathOperator*{\esssup}{ess\,sup}
\DeclareMathOperator*{\essinf}{ess\,inf}
\DeclareMathOperator*{\argmax}{arg\,max}
\DeclareMathOperator{\supp}{supp}
\DeclareMathOperator{\Var}{Var}
\DeclareMathOperator{\Cov}{Cov}
\DeclareMathOperator{\RV}{RV}
\DeclareMathOperator{\TV}{TV}
\newcommand{\Pp}{\mathbb P}
\newcommand{\E}{\mathbb E}
\newcommand{\R}{\mathbb R}
\newcommand{\1}{\mathbf 1}
\newcommand{\I}{I_\varepsilon}
\newcommand{\D}{\mathcal D_U}
\newcommand{\gU}{\gamma^\star_U}

%% --- présentation ---
\definecolor{gris}{gray}{0.45}
\definecolor{bleu}{RGB}{20,70,130}

\newcommand{\statut}[1]{%
  \nopagebreak\par\nopagebreak
  \hfill{\normalfont\footnotesize\textcolor{gris}{statut~: #1}}%
  \par\nopagebreak\vspace{-2pt}}

\newenvironment{idee}
  {\par\medskip\noindent\textcolor{bleu}{\textbf{L'idée.}}\ \itshape}
  {\par\medskip}

\newenvironment{apport}
  {\par\medskip\noindent\begingroup\color{bleu}\rule{\linewidth}{0.6pt}\par
   \nobreak\vspace{2pt}\noindent\textbf{Ce que cela apporte à l'article.}\ \normalcolor}
  {\par\vspace{2pt}\noindent\textcolor{bleu}{\rule{\linewidth}{0.6pt}}\endgroup\par\medskip}

\title{\bfseries Dossier de preuves\\[4pt]
\large Renforcement mathématique du criblage par indice de queue\\
\large en grande dimension}
\author{L.~Pommeret}
\date{\today}

\begin{document}
\maketitle

\begin{abstract}
\noindent
Ce dossier rassemble, sous forme démontrée et pédagogique, les résultats
obtenus au cours d'une campagne de renforcement du manuscrit
\emph{Tail-Index Sure Independence Screening for High-Dimensional
Heavy-Tailed Regression}. Chaque résultat est présenté selon le même
schéma~: l'énoncé, l'idée en langage ordinaire, la démonstration, puis
un encadré indiquant ce qu'il apporte à l'article --- ce qu'il remplace,
ce qu'il supprime, ou ce qu'il rend possible.

\smallskip\noindent
Trois acquis dominent. D'abord, \textbf{le facteur de largeur de bande
disparaît de l'échelle stochastique du score}~: celle-ci est $n\alpha$ et
non $n\alpha h$, ce qui améliore directement la condition de dimension
admissible du papier. Ensuite, \textbf{l'annexe des hypothèses primitives
peut être réellement primitive} sur une sous-classe qui contient les
quatre modèles de simulation, alors qu'on la croyait condamnée à porter
des conditions énoncées sur les objets projetés. Enfin, et c'est le
résultat central, \textbf{le score effectivement implémenté admet une
représentation linéaire par ligne}, ce qui autorise une valeur critique
calibrée sur le maximum --- sans changer d'estimateur, sans scission
d'échantillon, et sur une région de réglages qui contient les neuf
réglages du manuscrit.

\smallskip\noindent
Le dossier signale aussi, avec la même précision, ce qui est
\emph{impossible}~: une barrière de biais logarithmique interdit tout
théorème à écarts polynomialement rétrécissants dans la classe actuelle.
\end{abstract}

\tableofcontents

\bigskip
\hrule
\bigskip

\section*{Avant-propos : comment lire ce dossier}
\addcontentsline{toc}{section}{Avant-propos : comment lire ce dossier}

\paragraph{La règle de la campagne.} Aucun résultat n'a été retenu sur sa
seule apparence. Chaque énoncé produit a été soumis à une lecture
adverse, dont l'objet exclusif était de le mettre en défaut. Ce dossier
enregistre donc autant de \emph{rétrécissements} que de gains~: plusieurs
résultats annoncés comme «~exacts~» ne le sont que
conditionnellement, et le dire fait partie du résultat. Trois valeurs
numériques ont par ailleurs été recalculées indépendamment et sont
signalées comme telles.

\paragraph{Étiquettes de statut.} Chaque résultat porte une étiquette.

\begin{itemize}[leftmargin=2.2em,itemsep=2pt]
\item[\textbf{A}] \emph{Acquis} --- preuve complète, audit adverse rendu,
réparations intégrées à l'énoncé ci-dessous.
\item[\textbf{A$^-$}] \emph{Acquis sous réserve} --- preuve complète,
audit rendu, mais une clause reste à écrire proprement.
\item[\textbf{V}] \emph{En vérification} --- preuve complète, audit
adverse lancé, verdict non rendu à la date de ce dossier.
\item[\textbf{S}] \emph{Schéma} --- l'architecture de la preuve est
établie et chaque étape est nommée, mais la rédaction complète reste à
faire.
\item[\textbf{I}] \emph{Impossibilité} --- résultat négatif démontré.
\end{itemize}

\noindent
Cette distinction est importante~: un résultat marqué \textbf{S} ne doit
pas entrer dans l'article sans rédaction complète, même si son énoncé
est juste.

\paragraph{Conventions de notation.} Trois dérives ont été relevées au
cours de la campagne et sont closes par les conventions suivantes, qui
sont respectées partout dans ce dossier.

\begin{center}
\begin{tabular}{@{}ll@{}}
\toprule
Symbole & Sens \emph{unique} \\
\midrule
$\gU$ & $\esssup\gamma(U)$, maximum sur le \emph{support observé}.
        Jamais le maximum sur le cube. \\
$\xi_j(u)$ & $\esssup\{\gamma(U)\mid U_j=u\}$, profil projeté
             (indice de queue de $Y\mid U_j=u$). \\
$\Psi_j$ & $|\I|^{-1}\int_{\I}\xi_j(u)\,du$, profil intégré tronqué. \\
$\Delta_j$ & $\gU-\Psi_j\ \ (\ge 0)$, écart tronqué et défini
             sur le support. \\
$\D$ & $\{j:\Delta_j>0\}$, \emph{cible identifiable} du criblage
       marginal. \\
$s_D$ & $|\D|$. Distinct de $s_A$, taille d'une représentation
        structurelle minimale. \\
$\Delta_{\min}$ & $\min_{j\in\D}\Delta_j$. \\
$\alpha,h$ & fraction de queue locale et largeur de bande de rang~;
             $\alpha=n^{-a}$, $h=n^{-b}/2$. \\
\bottomrule
\end{tabular}
\end{center}

\noindent
La distinction entre $\D$ et l'ensemble structurel $A$ n'est pas
cosmétique~: c'est elle qui rend le problème identifiable. Le
chapitre~\ref{ch:cible} l'établit.

\paragraph{Ce que ce dossier ne fait pas.} Il ne réécrit pas l'article.
Il fournit la matière démontrée, et le chapitre~\ref{ch:changements}
indique section par section ce qu'il faudrait y changer. La rédaction
finale, en anglais et dans le style du manuscrit, reste à faire.

\clearpage

\section{La cible de population}\label{ch:cible}

Ce chapitre répond à une question que le manuscrit traitait trop vite~:
\emph{que vise exactement un criblage marginal par indice de queue~?}
La réponse n'est pas «~l'ensemble des variables dont dépend $\gamma$~».
Elle ne peut pas l'être, et le dire précisément est le premier acquis de
la campagne.

\subsection{L'identité de projection}

Tout repose sur un fait géométrique~: après projection sur une
coordonnée, la loi conditionnelle est un \emph{mélange}, et la queue d'un
mélange est gouvernée par sa composante la plus lourde. Le profil projeté
est donc une \emph{enveloppe supérieure}, pas une moyenne. Encore
faut-il le démontrer, et sous quelles hypothèses exactement.

Fixons le cadre. Soit $U$ à valeurs dans un espace borélien standard
$\mathcal X$, $Y>0$, et
\[
 \bar F_x(y)=\Pp(Y>y\mid U=x)
\]
une version régulière de la survie conditionnelle. Pour $c>1$, posons
\[
 \rho_c(g)=
 \begin{cases}
  c^{-1/g}, & g>0,\\
  0, & g=0,
 \end{cases}
\]
de sorte que la condition de variation régulière conditionnelle s'écrit
d'un seul tenant, indices nuls compris~:
\begin{equation}
 \frac{\bar F_x(cy)}{\bar F_x(y)}\longrightarrow\rho_c\{\gamma(x)\},
 \qquad c>1.
 \label{eq:rvc}
\end{equation}
La valeur $\rho_c(0)=0$ correspond à une queue à variation rapide.

Soit $K$ une loi de mélange quelconque sur $\mathcal X$ --- ce sera tour à
tour la loi de $U$ (queue inconditionnelle), la loi conditionnelle
$K_j(u,\cdot)$ de $U$ sachant $U_j=u$ (fibre), ou la loi conditionnelle à
$U_j$ dans une fenêtre de rangs. Écrivons
$\bar F_K(y)=\int\bar F_x(y)\,K(dx)$.

\begin{definition}[Loi inclinée par excès]
Lorsque $\bar F_K(y)>0$, la \emph{loi de mélange inclinée par excès} est
\begin{equation}
 \Pi_y(dx)=\frac{\bar F_x(y)\,K(dx)}{\bar F_K(y)}.
 \label{eq:tilt}
\end{equation}
\end{definition}

\noindent
C'est la loi de $U$ \emph{sachant que l'on observe un dépassement du
niveau $y$}. Elle contient tout le problème, en raison de l'identité
élémentaire
\begin{equation}
 \frac{\bar F_K(cy)}{\bar F_K(y)}
 =\int\frac{\bar F_x(cy)}{\bar F_x(y)}\,\Pi_y(dx).
 \label{eq:identite}
\end{equation}

\begin{idee}
La convergence ponctuelle \eqref{eq:rvc} ne suffit pas à passer à la
limite dans \eqref{eq:identite}, parce que la mesure d'intégration
$\Pi_y$ \emph{bouge avec $y$}. C'est là, et nulle part ailleurs, que se
niche toute la difficulté de la projection. Il faut donc deux clauses~:
l'une qui empêche les dépassements de se concentrer là où la convergence
n'a pas encore eu lieu, l'autre qui garantit qu'ils se concentrent bien
sur les indices maximaux.
\end{idee}

\begin{theoreme}[Projection mesurable sous uniformité inclinée]
\label{thm:projection}
\statut{A}
Soit $\gamma:\mathcal X\to[0,\Gamma]$ mesurable et
$m=\esssup_{x\sim K}\gamma(x)$. Supposons que pour tout $c>1$
\begin{equation}
 \int\Bigl|\frac{\bar F_x(cy)}{\bar F_x(y)}-\rho_c\{\gamma(x)\}\Bigr|
 \,\Pi_y(dx)\longrightarrow 0,
 \tag{TU1}\label{eq:TU1}
\end{equation}
et, si $m>0$, que pour tout $\eta>0$
\begin{equation}
 \Pi_y\{\gamma\le m-\eta\}\longrightarrow 0 .
 \tag{TU2}\label{eq:TU2}
\end{equation}
Alors
\[
 \frac{\bar F_K(cy)}{\bar F_K(y)}\longrightarrow c^{-1/m},
 \qquad c>1,
\]
c'est-à-dire que le mélange a pour indice de queue $m$. Si $m=0$, alors
$\gamma=0$ $K$-presque sûrement et \eqref{eq:TU1} seule donne
$\bar F_K(cy)/\bar F_K(y)\to0$.

Aucune continuité de $x\mapsto\gamma(x)$, aucune atteinte du supremum,
aucune borne inférieure strictement positive sur $\gamma$ ne sont
requises.
\end{theoreme}

\begin{proof}
Par \eqref{eq:identite},
\[
 \Bigl|\frac{\bar F_K(cy)}{\bar F_K(y)}-\rho_c(m)\Bigr|
 \le
 \underbrace{\int\Bigl|\frac{\bar F_x(cy)}{\bar F_x(y)}
   -\rho_c\{\gamma(x)\}\Bigr|\,\Pi_y(dx)}_{(\mathrm I)}
 +\underbrace{\Bigl|\int\rho_c\{\gamma(x)\}\,\Pi_y(dx)-\rho_c(m)
   \Bigr|}_{(\mathrm{II})}.
\]
Le terme $(\mathrm I)$ tend vers zéro par \eqref{eq:TU1}.

Si $m>0$~: la fonction $g\mapsto\rho_c(g)=c^{-1/g}$ est bornée et
continue sur $[0,\Gamma]$, et \eqref{eq:TU2} dit exactement que $\gamma$
converge vers $m$ en $\Pi_y$-probabilité. Le théorème de convergence
dominée pour la convergence en probabilité donne
$(\mathrm{II})\to0$.

Si $m=0$~: comme $\gamma\ge0$ et $\esssup_K\gamma=0$, on a $\gamma=0$
$K$-presque sûrement, donc $\rho_c\{\gamma(x)\}=0$ pour $K$-presque tout
$x$, et $\Pi_y\ll K$ entraîne que $(\mathrm{II})$ est identiquement
nul.
\end{proof}

\begin{remarque}[Ce qui est exactement nécessaire]
\label{rem:exactitude-projection}
La condition littéralement nécessaire et suffisante est
$\int \bar F_x(cy)/\bar F_x(y)\,\Pi_y(dx)\to\rho_c(m)$, puisque le membre
de gauche \emph{est} le rapport de queue projeté. La paire
\eqref{eq:TU1}--\eqref{eq:TU2} est suffisante, et \eqref{eq:TU2} est
nécessaire \emph{conditionnellement} à \eqref{eq:TU1}~; mais la paire
entière n'est pas nécessaire. Un audit a produit une loi projetée
exactement de Pareto qui viole l'uniformité inclinée $L^1$. Il ne faut
donc pas écrire que ces deux clauses caractérisent la projection.
\end{remarque}

\begin{apport}
Le manuscrit démontrait l'identité de projection sous continuité de
$\gamma$, support de fibre plein et non-atomicité. Le
théorème~\ref{thm:projection} supprime les trois. Il autorise en outre
les indices nuls dans le même énoncé, ce que la formulation en variation
régulière ordinaire ne permettait pas. Concrètement~: l'hypothèse (S)
telle qu'imprimée peut être retirée de cette étape, et la
remarque~\ref{rem:exactitude-projection} doit être ajoutée pour ne pas
surénoncer.
\end{apport}

\subsection{Le mélange Pareto exact : ni continuité, ni indices minorés}

Le cas Pareto exact mérite un énoncé séparé, parce qu'il est celui où la
suppression des hypothèses est la plus spectaculaire, et parce que sa
preuve --- une inclinaison exponentielle suivie d'un argument de
variation lente --- est le prototype de tout ce qui suit.

\begin{theoreme}[Projection Pareto exacte, indices nuls autorisés]
\label{thm:pareto}
\statut{A}
Supposons que pour $y\ge1$
\[
 \bar F_x(y)=
 \begin{cases}
  y^{-1/\gamma(x)}, & \gamma(x)>0,\\
  0, & \gamma(x)=0,
 \end{cases}
\]
et soit $m=\esssup_K\gamma$. Si $m>0$, alors
\[
 \bar F_K(y)=y^{-1/m}L_K(y),
\]
où $L_K$ est à variation lente. Le mélange a donc l'indice $m$, sans
continuité de $\gamma$, sans que $m$ soit atteint, et sans borne
inférieure strictement positive sur $\gamma$.
\end{theoreme}

\begin{idee}
On soustrait l'indice maximal en exposant. Il reste une transformée de
Laplace $B(t)=\E_K e^{-tZ}$ d'une variable positive $Z=1/\gamma-1/m$
dont l'\emph{infimum essentiel est nul}. Une transformée de Laplace en un
tel point est à variation lente~: c'est tout le contenu du théorème.
\end{idee}

\begin{proof}
Avec la convention $1/0=+\infty$, posons
$Z(x)=1/\gamma(x)-1/m\in[0,\infty]$. Alors
\[
 \bar F_K(y)=y^{-1/m}\,\E_K\bigl[e^{-(\log y)Z}\bigr],
\]
et l'on pose $B(t)=\E_K(e^{-tZ})$, $L_K(y)=B(\log y)$.

Comme $m$ est le supremum essentiel de $\gamma$, l'infimum essentiel de
$Z$ est nul. Introduisons la loi inclinée
\[
 \Pi^Z_t(A)=\frac{\E[e^{-tZ}\1\{Z\in A\}]}{B(t)} .
\]
Pour tout $\eta>0$,
\[
 \Pi^Z_t(Z>\eta)
 \le\frac{e^{-t\eta}}{\E[e^{-tZ}\1\{Z\le\eta/2\}]}
 \le\frac{e^{-t\eta}}{e^{-t\eta/2}\Pp(Z\le\eta/2)}
 =\frac{e^{-t\eta/2}}{\Pp(Z\le\eta/2)}\longrightarrow0,
\]
le dénominateur étant strictement positif précisément parce que
$\essinf Z=0$. Donc $Z\to0$ en $\Pi^Z_t$-probabilité. Il vient, pour
tout $a>0$,
\[
 \frac{B(t+a)}{B(t)}=\E_{\Pi^Z_t}(e^{-aZ})\longrightarrow1,
\]
par convergence dominée ($0\le e^{-aZ}\le1$). Ainsi
$L_K(cy)/L_K(y)=B(\log y+\log c)/B(\log y)\to1$ pour $c>1$, puis pour
$0<c<1$ par inversion. Donc $L_K$ est à variation lente.
\end{proof}

\noindent
Le théorème couvre~: $\gamma$ discontinue~; $\gamma=0$ sur un ensemble de
mesure positive arbitraire~; $\gamma(x)\downarrow0$ au bord~; et un
supremum approché mais jamais atteint. La partie d'indice nul ne
contribue rien à la queue asymptotique dès que $m>0$.

\subsection{Pourquoi la cible doit être définie sur le support}

Le manuscrit écrivait, sous continuité de $\gamma$,
$\gamma^\star=\max_{x\in\supp(P_U)}\gamma(x)$. Cette écriture n'est pas
seulement inconfortable pour $\gamma$ mesurable~: elle est
\emph{illégitime}, car ce n'est pas une fonctionnelle de la loi.

\begin{proposition}[Le maximum sur le support n'est pas une fonctionnelle de la loi]
\label{prop:pas-fonctionnelle}
\statut{A}
Soit $U\sim\mathrm{Unif}(0,1)$. Considérons deux versions de la fonction
d'indice qui coïncident pour tout $u\ne0$ mais diffèrent en $u=0$~:
\[
 \gamma_0(u)=0,\qquad \gamma_1(u)=\1\{u=0\}.
\]
Les lois jointes de $(U,Y)$ engendrées sont \emph{identiques}, alors que
\[
 \max_{u\in[0,1]}\gamma_0(u)=0,
 \qquad
 \max_{u\in[0,1]}\gamma_1(u)=1 .
\]
En revanche les deux suprema essentiels valent $0$.
\end{proposition}

\begin{proof}
Les deux versions diffèrent sur $\{0\}$, de mesure de Lebesgue nulle~;
une version de la loi conditionnelle n'est définie qu'à un ensemble
négligeable près, donc les lois jointes coïncident. Les deux maxima
ponctuels se lisent directement.
\end{proof}

\noindent
La seule définition invariante par changement de version est donc
\begin{equation}
 \boxed{\ \gU=\esssup_{P_U}\gamma(U).\ }
 \label{eq:cible}
\end{equation}

\begin{apport}
Il faut remplacer partout dans l'article $\max_{\supp}\gamma$ par
$\esssup\gamma$, y compris là où la continuité rendait l'égalité vraie~:
l'égalité est vraie, mais l'écriture par le maximum ne l'est pas si
$\gamma$ n'est connue qu'à une version près, et c'est le cas dès qu'on ne
suppose plus la continuité. C'est une correction de fond, pas de style.
\end{apport}

\subsection{L'enveloppe restreinte au support et la cible exacte}

Sous continuité, on peut aller plus loin et caractériser exactement
quelles coordonnées un criblage marginal peut voir. Soit
$K_j(u,\cdot)$ une version ponctuelle de la loi de $U_{-j}$ sachant
$U_j=u$, de support $C_j(u)\subseteq[0,1]^{p-1}$ (non vide, compact), et
soit $\mathcal M=\argmax_{[0,1]^p}\gamma$. Définissons l'\emph{ensemble
des argmax accessibles}
\[
 \mathcal H_j=\bigl\{u\in\I:(\{u\}\times C_j(u))\cap\mathcal M\ne\varnothing\bigr\}.
\]

\begin{theoreme}[Enveloppe restreinte au support et cible marginale exacte]
\label{thm:enveloppe}
\statut{A}
Sous les hypothèses de variation régulière conditionnelle et $\gamma$
continue --- \emph{sans aucune hypothèse de support plein} --- on a pour
toute version ponctuelle
\begin{equation}
 \xi_j(u)=\esssup_{v\sim K_j(u,\cdot)}\gamma\{\iota_j(u,v)\}
        =\max_{v\in C_j(u)}\gamma\{\iota_j(u,v)\},
 \label{eq:env}
\end{equation}
et par conséquent
\[
 \xi_j(u)=\gamma^\star\iff u\in\mathcal H_j,
 \qquad
 \Delta_j>0\iff\lambda(\I\setminus\mathcal H_j)>0 .
\]
La dernière équivalence ne requiert aucune continuité de
$u\mapsto\xi_j(u)$.
\end{theoreme}

\begin{proof}
Le théorème de projection donne la première expression de \eqref{eq:env}.
Fixons $j,u$ et abrégeons $K=K_j(u,\cdot)$, $C=C_j(u)$,
$f(v)=\gamma\{\iota_j(u,v)\}$. Comme $C$ est compact et $f$ continue,
$m=\max_Cf$ est atteint. Puisque $K(C)=1$, on a $f\le m$
$K$-p.s., donc $\esssup_Kf\le m$. Soit $v_\star\in C$ atteignant $m$~:
pour tout $\eta>0$, la continuité fournit un voisinage ouvert $O_\eta$ de
$v_\star$ sur lequel $f>m-\eta$, et la propriété caractéristique du
support donne $K(O_\eta)>0$, d'où $\esssup_Kf\ge m-\eta$. On conclut en
faisant $\eta\downarrow0$.

Le maximum dans \eqref{eq:env} vaut $\gamma^\star$ exactement lorsqu'il
est atteint en un point de l'argmax global, ce qui est la deuxième
assertion. Enfin $0\le\gamma^\star-\xi_j(u)$ et
$\Delta_j=|\I|^{-1}\int_{\I}\{\gamma^\star-\xi_j(u)\}\,du$~; une fonction
mesurable positive est d'intégrale nulle si et seulement si elle est
nulle presque partout.
\end{proof}

\begin{corollaire}[Identification minimale]
\label{cor:identification}
\statut{A}
La condition d'\emph{accessibilité nulle}
\[
 \lambda(\I\setminus\mathcal H_j)=0\quad\text{pour tout }j\notin A
 \tag{NA}
\]
équivaut à $\D\subseteq A$. On a $\D=A$ si et seulement si, en outre,
$\lambda(\I\setminus\mathcal H_j)>0$ pour tout $j\in A$. Le support de
fibre plein est \emph{suffisant} pour ces conclusions, mais nullement
nécessaire.
\end{corollaire}

\begin{remarque}[Pourquoi une variable inactive peut porter un vrai signal]
Si (NA) échoue, une coordonnée inactive peut restreindre le support
conditionnel des coordonnées actives de manière à supprimer les
configurations presque maximisantes sur un ensemble de valeurs de mesure
positive. Son indice projeté varie alors, \emph{bien que} $\gamma$ ne
dépende pas d'elle. Ce n'est pas une erreur d'estimation~: la loi
marginale de $(U_j,Y)$ contient réellement un signal d'indice de queue.
Aucune méthode marginale coordonnée par coordonnée ne peut distinguer un
tel relais d'une variable structurellement active sans hypothèse
d'accessibilité ou d'indépendance conditionnelle.
\end{remarque}

\begin{apport}
C'est le résultat qui rend le problème honnête. L'article visait $A$,
l'ensemble structurel~; il doit viser $\D$, et présenter
$\D=A$ comme une \emph{hypothèse d'accessibilité} explicite, non comme un
acquis. Cela remplace toute «~condition de force de signal générique~» par
une caractérisation exacte de ce qu'un criblage marginal peut voir, ce qui
est à la fois plus faible en hypothèses et plus fort en conclusion.
\end{apport}

\subsection{Le cas sans maximiseur, et le cas d'indice nul}

Deux dégénérescences doivent être traitées, car elles sont réelles.

\paragraph{Pas de maximiseur.} Si $\gamma$ est seulement mesurable,
l'ensemble $M=\{x:\gamma(x)=\gU\}$ peut être vide ou négligeable alors
même que $\gU>0$. La caractérisation par contact avec l'argmax tombe. On
la remplace par une version «~presque maximale~»~: pour $\eta>0$, soit
$M_\eta=\{x:\gamma(x)>\gU-\eta\}$.

\begin{theoreme}[Caractérisation mesurable par écart positif]
\label{thm:gap-mesurable}
\statut{A}
Pour presque tout $u$,
\[
 \xi_j(u)=\gU
 \iff
 K_j(u,M_\eta)>0\ \text{ pour tout }\eta>0,
\]
et
\[
 \Delta_j>0
 \iff
 \exists\,\eta>0:\ \lambda\{u\in\I:K_j(u,M_\eta)=0\}>0 .
\]
\end{theoreme}

\begin{proof}
Si $K_j(u,M_\eta)>0$ pour tout $\eta>0$, alors
$\xi_j(u)\ge\gU-\eta$ pour tout $\eta$, donc $\xi_j(u)\ge\gU$~;
l'inégalité inverse est toujours vraie, d'où l'égalité. Réciproquement,
si $K_j(u,M_\eta)=0$ pour un $\eta$, alors $\xi_j(u)\le\gU-\eta$.

Pour la seconde équivalence~: si le membre de droite vaut, alors
$\gU-\xi_j(u)\ge\eta$ sur un ensemble de mesure positive, donc
$\Delta_j>0$. Réciproquement, si $\Delta_j>0$ alors
$\lambda\{\xi_j<\gU\}>0$~; comme
$\{\xi_j<\gU\}=\bigcup_{m\ge1}\{\xi_j\le\gU-1/m\}$, il existe $\delta>0$
tel que $A_\delta=\{u:\xi_j(u)\le\gU-2\delta\}$ soit de mesure positive,
et pour $u\in A_\delta$ on a $K_j(u,M_\delta)=0$.
\end{proof}

\begin{exemple}[L'ancienne formulation par contact échoue]
Soit $U=(U_1,U_2)$ uniforme sur le carré et
$\gamma(u_1,u_2)=u_2$ pour $u_2<1$, $\gamma(u_1,1)=0$. Alors $\gU=1$,
mais $\gamma$ n'atteint jamais $1$, donc $M=\varnothing$. Pourtant
$\xi_1(u_1)=\esssup_{u_2}u_2=1$ pour tout $u_1$, et $\Delta_1=0$. La
formulation presque maximale est correcte~: toute fibre verticale a une
masse strictement positive dans chaque $M_\eta$.
\end{exemple}

\paragraph{Indice nul.} Puisque $\gamma\ge0$,
\[
 \gU=0\ \Longrightarrow\ \gamma(U)=0\ \ P_U\text{-p.s.}
 \ \Longrightarrow\ \xi_j\equiv0,\ \Psi_j=0,\ \Delta_j=0,\ \D=\varnothing .
\]

\begin{exemple}[La cible est aveugle aux queues légères]
\label{ex:aveugle}
Soit $p=1$, $U\sim\mathrm{Unif}(0,1)$ et
\[
 \Pp(Y>y\mid U=u)=
 \begin{cases}
  e^{-y}, & u\le1/2,\\
  e^{-y^2}, & u>1/2 .
 \end{cases}
\]
Les deux lois conditionnelles ont l'indice zéro. Leurs queues sont
pourtant très différentes, mais
$\gU=\xi_1(u)=\Delta_1=0$.
\end{exemple}

\begin{apport}
Ce n'est pas un échec d'estimation, c'est une \emph{limite de la cible},
et l'article doit l'énoncer. Conserver le score de Hill est une
restriction assumée~: la théorie couvre les indices de Fréchet positifs,
tandis que les procédures de type moment conditionnel couvrent Fréchet,
Gumbel et Weibull. Puisque $\gU=0\Rightarrow\D=\varnothing$, la portée du
papier doit être annoncée comme telle en introduction, plutôt que
découverte par le lecteur.
\end{apport}

\subsection{Trois impossibilités, et leur portée exacte}

La campagne a produit trois résultats négatifs. Ils sont utiles --- ils
justifient des hypothèses au lieu de les subir --- à condition d'énoncer
leur portée avec exactitude, car la première rédaction les avait tous
trois surénoncés.

\begin{theoreme}[Apparition retardée]
\label{thm:retard}
\statut{I, A}
Il existe des suites de modèles, satisfaisant la variation régulière
conditionnelle en chaque point, pour lesquelles aucune procédure ne peut
séparer les coordonnées actives des inactives uniformément sur des
tableaux triangulaires, l'onset de la régularité pouvant être retardé
arbitrairement.
\end{theoreme}

\noindent
\emph{Portée.} L'impossibilité est \textbf{uniforme sur tableaux
triangulaires}, et non ponctuelle à loi fixée. Il faut écrire «~pas
détectable uniformément~». Elle montre qu'\emph{une} restriction excluant
l'apparition arbitrairement retardée est nécessaire --- \textbf{pas} que
la condition (P2) du manuscrit soit \emph{la} condition nécessaire. Enfin
$1/\log n$ et $1/\log(1/\alpha_n)$ ne coïncident que si
$\log(1/\alpha_n)\asymp\log n$.

Un témoin élémentaire du mécanisme~: soit $U\sim\mathrm{Unif}(0,1)$,
$0<\beta<1$, et conditionnellement à $U=u$,
$Y=P$ avec probabilité $1/2$ et $Y=u^{-1/\beta}$ avec probabilité $1/2$,
où $P$ est Pareto standard indépendant. Pour chaque $u>0$,
$\Pp(Y>y\mid U=u)\sim\tfrac12y^{-1}$, donc \emph{toute} loi
conditionnelle a l'indice $1$. Mais
\[
 \Pp(Y>y)=\tfrac12y^{-1}+\tfrac12 y^{-\beta}\sim\tfrac12y^{-\beta},
\]
et le mélange a l'indice $1/\beta>1$. La variation régulière ponctuelle
ne se transmet pas au mélange sans contrôle incliné.

\begin{theoreme}[Invisibilité marginale]
\label{thm:invisible}
\statut{I, A}
Il existe deux modèles d'ensembles actifs différents dont toutes les lois
marginales bivariées $(U_j,Y)$ coïncident.
\end{theoreme}

\noindent
\emph{Portée.} L'énoncé ne vaut que pour une décision mesurable par
rapport au \textbf{seul couple} $(U_j,Y)$. Une procédure regardant une
autre coordonnée distingue les modèles. Il est faux d'écrire que «~toute
règle marginale a une erreur $\ge1/2$~». L'énoncé fort correct porte sur
la \emph{liste complète} des marginales bivariées, et se réalise avec
$p=3$ et des phases décalées.

\begin{theoreme}[Non-identifiabilité sous dépendance libre]
\label{thm:nonid}
\statut{I, A}
Sous dépendance non contrainte des covariables, $A$ n'est pas une
fonctionnelle de la loi jointe.
\end{theoreme}

\noindent
\emph{Portée.} Confirmée, et \emph{plus forte} que la version d'origine.
La réparation consiste à viser $\D$, ou à exiger
$\supp(U)=[0,1]^p$ \textbf{et} $A=\D$. L'hypothèse (S) telle
qu'imprimée dans le papier ne suffit pas, non plus que (NS).

\begin{apport}
Ces trois énoncés remplacent, dans la discussion du papier, des phrases
prudentes du type «~un criblage marginal peut échouer~» par des théorèmes.
Leur rôle éditorial est de \emph{justifier} les hypothèses
d'accessibilité et de régularité~: on montre qu'aucune hypothèse de ce
type ne peut être supprimée, plutôt que de les poser sans commentaire.
Il faut en revanche corriger les trois portées comme ci-dessus~: en
l'état, le manuscrit les énoncerait trop fort, ce qu'un rapporteur
attentif verrait.
\end{apport}

\clearpage

\section{L'annulation de la largeur de bande}\label{ch:h}

C'est le premier gain quantitatif de la campagne, et le plus directement
lisible dans les énoncés du papier. Il a été obtenu, puis reconfirmé par
trois routes indépendantes~: le coloriage présenté ici, un score par blocs
disjoints, et un calcul par fonction d'influence.

\subsection{Le problème}

Le manuscrit énonce que l'erreur maximale du score sur $p$ coordonnées est
d'ordre $\sqrt{\log(pn)/(n\alpha h)}$, où $n\alpha h$ est le nombre
d'extrêmes locaux dans une fenêtre. Cette écriture est naturelle~: chaque
profil local $\widehat\xi_j(u)$ repose sur environ $n\alpha h$ excès, donc
fluctue à l'échelle $(n\alpha h)^{-1/2}$.

\begin{idee}
Mais le score n'est pas un profil local~: c'est une \emph{moyenne} de
profils locaux sur toute la grille de centres. Or on ne moyenne pas des
erreurs indépendantes de la même façon qu'on borne la pire d'entre elles.
Les fenêtres se recouvrent, donc les profils ne sont pas indépendants~;
mais deux fenêtres dont les centres sont assez éloignés sont
\emph{disjointes}, donc indépendantes conditionnellement à la colonne de
covariables. Il suffit alors de colorier la grille de manière que chaque
classe de couleur ne contienne que des fenêtres disjointes. Le nombre de
couleurs nécessaire est exactement la largeur d'une fenêtre --- et c'est
ce facteur qui, en se compensant avec la taille de la grille, fait
disparaître $h$.
\end{idee}

\noindent
Le point crucial de rédaction~: il faut moyenner l'expansion
\emph{signée}, puis prendre la valeur absolue. Le manuscrit prenait les
valeurs absolues avant de moyenner, ce qui détruit toute compensation.

\subsection{Mise en place}

Écrivons
\[
 \mathcal R_n=\{r\le n:\ r/(n+1)\in\I\},\quad L_n=|\mathcal R_n|,\quad
 d_n=\lfloor h(n+1)\rfloor .
\]
Pour un centre $r\in\mathcal R_n$, la fenêtre de rangs est
$W_r=\{q:|q-r|\le d_n\}$, de cardinal $m_r$, et l'on garde
$k_r=\lfloor\alpha m_r\rfloor$ excès. Posons
$K_n=\min_r k_r$ et
\[
 \chi_n=2d_n+1,\qquad
 v_n=\frac{\chi_n}{L_n^2}\sum_{r\in\mathcal R_n}\frac1{k_r},\qquad
 c_n=\frac{\chi_n}{L_nK_n},\qquad
 \rho_n(x)=\sqrt{2v_nx}+c_nx .
\]

\noindent
La quantité $\chi_n$ est le \emph{nombre chromatique} du graphe de
recouvrement des fenêtres~: deux centres de même résidu modulo $\chi_n$
diffèrent d'au moins $2d_n+1$, donc leurs fenêtres sont disjointes.

\subsection{Les deux lemmes}

\begin{lemme}[Fonction génératrice des moyennes Gamma]
\label{lem:gamma}
\statut{A}
Soit $G_k=k^{-1}\sum_{i\le k}E_i$ avec $E_i$ i.i.d.\ exponentielles
standard. Alors pour $|\lambda|<k$,
\[
 \log\E\exp\{\lambda(G_k-1)\}\le\frac{\lambda^2}{2k\{1-|\lambda|/k\}} .
\]
\end{lemme}

\begin{proof}
Pour $0\le\lambda<k$, avec $x=\lambda/k$,
\[
 \log\E e^{\lambda(G_k-1)}=k\{-x-\log(1-x)\}\le\frac{kx^2}{2(1-x)},
\]
la dernière inégalité s'obtenant en intégrant $x/(1-x)\le x/(1-x_0)$ sur
$[0,x_0]$, ou directement par la série entière. Pour $-k<\lambda<0$, on
applique le même argument à $x=-\lambda/k$ et l'on utilise
$x-\log(1+x)\le x^2/2$.
\end{proof}

\begin{lemme}[Concentration d'une moyenne d'espacements locaux recouvrants]
\label{lem:coloriage}
\statut{A}
Fixons une coordonnée $j$ et conditionnons sur sa colonne de covariables.
Pour $r\in\mathcal R_n$, soit $G_{jr}$ la variable de Rényi apparaissant
dans le développement local de Hill à position gelée~; $G_{jr}$ suit donc
la loi d'une moyenne de $k_r$ exponentielles standard indépendantes. Si
$|a_r|\le A$ et
\[
 S_a=\frac1{L_n}\sum_{r\in\mathcal R_n}a_r(G_{jr}-1),
\]
alors pour tout $x>0$,
\[
 \Pp\{|S_a|>A\rho_n(x)\mid\bm U_j\}\le2e^{-x} .
\]
\end{lemme}

\begin{proof}
Partitionnons $\mathcal R_n$ selon les résidus modulo $\chi_n=2d_n+1$.
Deux centres d'une même classe diffèrent d'au moins $2d_n+1$, donc leurs
fenêtres de rangs sont disjointes. Conditionnellement à la colonne de
covariables, les variables issues de fenêtres disjointes sont
indépendantes~; les $G_{jr}$ d'une même classe de couleur le sont donc.

Soient $S_1,\dots,S_{\chi_n}$ les sommes par couleur. L'inégalité de
Hölder et le lemme~\ref{lem:gamma} donnent, dès que
$A\chi_n|\theta|/(L_nK_n)<1$,
\[
 \log\E(e^{\theta S_a}\mid\bm U_j)
 \le\frac1{\chi_n}\sum_{c=1}^{\chi_n}\log\E(e^{\chi_n\theta S_c}\mid\bm U_j)
 \le\frac{A^2\theta^2v_n}{2\{1-Ac_n|\theta|\}} .
\]
Ainsi $S_a$ est sous-gamma bilatère de facteur de variance $A^2v_n$ et de
facteur d'échelle $Ac_n$~; l'optimisation de Chernoff donne la borne.
\end{proof}

\begin{remarque}[Où passe le facteur $h$]
Le coloriage coûte un facteur $\chi_n\asymp nh$ dans la borne de Hölder,
mais le nombre de termes moyennés est $L_n\asymp n$ et chacun apporte
$1/k_r\asymp1/(n\alpha h)$. Le produit
$\chi_n\cdot\frac1{L_n}\cdot\frac1{n\alpha h}\asymp\frac{1}{n\alpha}$~:
les deux $h$ se compensent \emph{exactement}. C'est toute l'affaire.
\end{remarque}

\subsection{Le théorème de concentration}

\begin{theoreme}[Concentration à échantillon fini du score glissant d'origine]
\label{thm:score}
\statut{A}
Sous les hypothèses du développement à position gelée du manuscrit, avec
$K_n\ge2$, $0<\alpha\le1/3$, $\eta>0$, $0<\tau\le3\alpha$ et $x>0$~:
\begin{align*}
 \Pp\Bigl\{\max_{j\le p}|\widehat\Psi_j-\Psi_j|
 &>q_n+2\omega_n(\eta,\tau)+b_n+(\Gamma+b_n)\rho_n(x)\Bigr\}\\
 &\le 2pe^{-2n\eta^2}+pn\tau+pL_ne^{-(K_n+1)/4}+4pe^{-x} .
\end{align*}
Le terme stochastique est gouverné par $v_n$ et $c_n$, et \emph{non} par
$K_n^{-1}$ seul.
\end{theoreme}

\begin{proof}
Fixons $j$ et un centre $r$. Sur l'événement de déplacement de rangs,
l'événement de troncature globale, et l'événement que la
$(k_r+1)$-ième statistique d'ordre de chaque fenêtre est au plus
$3\alpha$, le calcul à position gelée s'écrit \emph{sous forme signée}
\[
 \widehat\xi_j(r/(n+1))-\xi_j(r/(n+1))
 =e_{jr}+\xi_j(r/(n+1))(G_{jr}-1)+R_{jr},
\]
avec $|e_{jr}|\le2\omega_n(\eta,\tau)$ et $|R_{jr}|\le b_nG_{jr}$. C'est
la décomposition du manuscrit~; le seul changement est qu'elle est
moyennée \emph{avant} passage aux valeurs absolues.

Posons
$A_j=L_n^{-1}\sum_r\xi_j(r/(n+1))(G_{jr}-1)$ et
$C_j=L_n^{-1}\sum_r(G_{jr}-1)$. En moyennant l'expansion signée puis en
appliquant le lemme de quadrature, on obtient trajectoire par
trajectoire
\[
 |\widehat\Psi_j-\Psi_j|\le q_n+2\omega_n(\eta,\tau)+|A_j|+b_n(1+|C_j|).
\]
Le lemme~\ref{lem:coloriage}, appliqué d'abord avec
$a_r=\xi_j(r/(n+1))$ puis avec $a_r=1$, donne
$\Pp\{|A_j|>\Gamma\rho_n(x)\mid\bm U_j\}\le2e^{-x}$ et
$\Pp\{|C_j|>\rho_n(x)\mid\bm U_j\}\le2e^{-x}$~; l'union sur $j$ contribue
$4pe^{-x}$.

L'événement de déplacement de rangs échoue avec probabilité au plus
$2pe^{-2n\eta^2}$, la troncature globale avec probabilité au plus
$pn\tau$, et, conditionnellement à une colonne,
$\Pp\{V_{(k_r+1)}>3\alpha\mid\bm U_j\}\le e^{-(k_r+1)/4}\le e^{-(K_n+1)/4}$,
d'où le troisième terme par union sur les $pL_n$ fenêtres.
\end{proof}

\begin{corollaire}[La largeur de bande disparaît du taux stochastique]
\label{cor:h}
\statut{A}
Si $nh\ge4$, $\alpha m_n^-\ge2$ et $L_n\ge c_\varepsilon n$, alors
\[
 v_n\vee c_n\le\frac{C_\varepsilon}{n\alpha},
\]
et donc, pourvu que $K_n\gg\log(pn)$,
\[
 \max_{j\le p}|\widehat\Psi_j-\Psi_j|
 =O_{\Pp}\!\left(\sqrt{\frac{\log p}{n\alpha}}+\frac{\log p}{n\alpha}\right)
 +q_n+2\omega_n(\eta_n,\tau_n)+b_n .
\]
Si le biais déterministe est $o(\Delta_{\min})$, le criblage sûr découle
de
\[
 \boxed{\ \log p=o(n\alpha\,\Delta_{\min}^2)\ }
\]
jointement aux conditions séparées de faisabilité locale et de
localisation.
\end{corollaire}

\begin{proof}
La géométrie déterministe des rangs donne $K_n\ge n\alpha h/4$, et
$\chi_n=2\lfloor h(n+1)\rfloor+1\le Cnh$ pour $n$ grand. Donc
\[
 v_n\le\frac{\chi_n}{L_nK_n}\le\frac{C_\varepsilon\,nh}{n\cdot n\alpha h/4}
 =\frac{C_\varepsilon'}{n\alpha},
\]
et de même pour $c_n$. On prend ensuite $x$ proportionnel à $\log p$ dans
le théorème~\ref{thm:score}, et l'on conclut par le lemme déterministe de
séparation du score.
\end{proof}

\subsection{Où $n\alpha h$ survit malgré tout}

Il serait faux d'en conclure que $h$ disparaît du papier. Il disparaît de
\emph{l'échelle stochastique du score intégré}. Il reste indispensable
dans trois rôles, et les confondre serait une erreur~:

\begin{enumerate}[leftmargin=2em,itemsep=3pt]
\item \textbf{Faisabilité locale}~: il faut $K_n\asymp n\alpha h\to\infty$,
et même $n\alpha h\gg\log(pn)$, pour que chaque seuil local tombe
effectivement dans la queue. C'est ce qui donne $a+b<1$.
\item \textbf{Biais spatial}~: la fenêtre a une largeur, donc le profil
qu'elle vise n'est pas $\xi_j(u)$ mais un supremum local. C'est l'objet du
chapitre~\ref{ch:spatial}.
\item \textbf{Concentration ponctuelle}~: si l'on veut une garantie sur
$\widehat\xi_j(u)$ \emph{en un point}, et non sur la moyenne, l'échelle
$(n\alpha h)^{-1/2}$ reste la bonne.
\end{enumerate}

\begin{apport}
C'est l'amélioration la plus directement visible. La condition de
criblage sûr du manuscrit,
$\log(pn)=o(n\alpha h\,\Delta_{\min}^2)$, devient
$\log p=o(n\alpha\,\Delta_{\min}^2)$~: on gagne un facteur $h^{-1}$, soit
$n^{b}$, sur la dimension admissible. Le prix est nul --- ni hypothèse
supplémentaire, ni changement d'estimateur~; il fallait seulement moyenner
avant de majorer. Concrètement, il faut réécrire l'abrégé, l'énoncé du
théorème de score, celui du théorème de criblage, et la discussion des
taux. Il faut aussi \emph{conserver} $n\alpha h$ aux trois endroits
ci-dessus, et le dire, faute de quoi la correction paraîtra trop belle.
\end{apport}

\clearpage

\section{La couche spatiale}\label{ch:spatial}

Le chapitre précédent a traité le hasard. Celui-ci traite la géométrie~:
que vise exactement une fenêtre de largeur non nulle, et à quelle
condition ce qu'elle vise est-il assez proche de ce qu'on veut~?

\subsection{Les fenêtres visent un supremum mobile, pas une moyenne}

C'est le point que le manuscrit passait sous silence. Sur une fenêtre
$W_h(u)=[u-h,u+h]\cap[0,1]$, la loi de $Y$ sachant $U_j\in W_h(u)$ est
encore un mélange. Par le théorème~\ref{thm:projection}, son indice est
donc le \emph{supremum} de $\xi_j$ sur la fenêtre, et non la moyenne de
$\xi_j$ sur la fenêtre~:
\[
 M_h\xi_j(u)=\sup_{v\in W_h(u)}\xi_j(v)\ \ge\ \xi_j(u).
\]
Le biais spatial est donc \emph{unilatéral et positif}. Une condition de
continuité du profil ne suffit pas à le contrôler~: c'est une condition
sur l'\emph{enveloppe supérieure} qu'il faut.

\begin{definition}[Enveloppe essentielle supérieure]
$\xi_j^\sharp(u)=\lim_{h\downarrow0}M_h\xi_j(u)$.
\end{definition}

\begin{proposition}[Condition spatiale exacte à loi fixe]
\label{prop:EUR}
\statut{A}
À loi fixe, la convergence du score de fenêtre de population vers
$\Psi_j$ quand $h\downarrow0$ équivaut \emph{exactement} à
\[
 \xi_j^\sharp=\xi_j\quad\text{presque partout.}
 \tag{EUR}
\]
\end{proposition}

\begin{remarque}[Un ensemble de Cantor gras comme borne de validité]
La condition (EUR) n'est pas automatique~: si $\xi_j$ est modifiée sur un
ensemble de Cantor gras (fermé, d'intérieur vide, de mesure positive),
l'enveloppe supérieure diffère du profil sur un ensemble de mesure
positive, et le score de fenêtre ne converge pas vers $\Psi_j$.
\end{remarque}

\subsection{Pourquoi (EUR) ne suffit pas en tableau triangulaire}

C'est le défaut le plus sérieux découvert par la campagne, et il portait
sur un résultat annoncé comme «~exact~».

\begin{theoreme}[Insuffisance de la forme qualitative]
\label{thm:damier}
\statut{I, A}
Il existe une suite de lois, indexée par $n$, telle que
$\xi^\sharp_{j,n}=\xi_{j,n}$ presque partout \emph{pour chaque $n$}, et
pourtant
\[
 \int_0^1 M_{h_n}f_n(u)\,du\longrightarrow1,
 \qquad
 \int_0^1 f_n(u)\,du\longrightarrow\tfrac12 .
\]
La condition (EUR) sous forme qualitative ne contrôle donc rien
uniformément lorsque la loi varie avec $n$.
\end{theoreme}

\begin{idee}
Le témoin est un \emph{damier de plus en plus fin}. À chaque $n$, le
profil alterne entre deux valeurs sur des intervalles de longueur bien
plus petite que $h_n$. En chaque point, l'enveloppe supérieure coïncide
presque partout avec le profil --- l'égalité (EUR) est donc satisfaite,
pour chaque $n$. Mais une fenêtre de largeur $h_n$ contient un grand
nombre d'oscillations complètes, donc son supremum vaut la valeur haute
partout, tandis que la moyenne du profil vaut la demi-somme. L'écart ne
tend pas vers zéro. Une égalité \emph{presque partout} est insensible à
la finesse~; c'est précisément ce que le régime triangulaire exploite.
\end{idee}

\noindent
La condition correcte est \emph{quantitative}~:
\begin{equation}
 \Omega_n(h_n)=\max_{j\le p_n}\frac1{|\I|}\int_{\I}
   \bigl\{M_{h_n}\xi_{j,n}(u)-\xi_{j,n}(u)\bigr\}\,du
 \longrightarrow0,
 \label{eq:omega}
\end{equation}
et, pour le criblage,
\begin{equation}
 \boxed{\ \Omega_n(h_n)=o(\Delta_{\min,n}).\ }
 \label{eq:omega-screen}
\end{equation}

\begin{apport}
Partout où le manuscrit --- ou une version antérieure de ce dossier ---
énonce que (EUR) est «~la condition spatiale exacte~», il faut écrire~:
(EUR) est exacte \emph{à loi fixe}~; en tableau triangulaire, c'est
\eqref{eq:omega-screen} qui vaut. C'est une correction obligatoire, car
le papier travaille en régime triangulaire ($p_n$ croît avec $n$).
\end{apport}

\subsection{Le rayon réellement utilisé par les fenêtres de rangs}

Une subtilité supplémentaire, et elle est agréable~: les fenêtres sont
construites sur des \emph{rangs empiriques}, donc leur extension en
échelle uniforme n'est pas $h_n$ mais un rayon aléatoire légèrement plus
grand.

\begin{proposition}[Rayon de travail]
\label{prop:rayon}
\statut{A}
La condition spatiale doit être imposée au rayon
\[
 r_n=h_n+2\sqrt{\frac{h_nx_n}{n}}+\frac{2x_n}{n},
 \qquad x_n\asymp\log(p_nn),
\]
et non à $h_n$. Avec ce rayon, l'appartenance à une fenêtre de rangs
n'introduit \emph{aucune} dépendance cachée, pourvu que la preuve soit
recentrée sur l'ancre aléatoire $U_{(r)j}$ plutôt que sur $r/(n+1)$, que
la condition soit posée sur un intervalle tronqué élargi, et que l'on
ajoute l'événement négligeable «~toutes les ancres restent dans cet
intervalle~».
\end{proposition}

\begin{idee}
Pour un centre de rang $r$ et une demi-largeur $d=\lfloor h(n+1)\rfloor$,
l'écart en échelle uniforme entre les extrémités de la fenêtre est
$U_{(r+d)}-U_{(r)}\sim\mathrm{Beta}(d,n+1-d)$. L'identité
bêta--binomiale fournit une borne exponentielle dont l'inversion donne
exactement le rayon ci-dessus~: le terme $\sqrt{h_nx_n/n}$ est la
fluctuation typique, le terme $x_n/n$ le régime de grandes déviations. Le
recentrage sur l'ancre \emph{aléatoire} est ce qui supprime l'ancien
déplacement global de type Dvoretzky--Kiefer--Wolfowitz.
\end{idee}

\subsection{Ce que coûte une largeur de bande dégénérée : le calcul explicite}

Le manuscrit observe, sans l'expliquer, que la colonne $b=0$ de sa grille
de réglages «~n'est jamais compétitive~». On peut maintenant le calculer.

Pour les modèles de simulation, $\gamma$ décroît le long de chaque
coordonnée active, de sorte que le profil projeté est explicitement
$\xi_j(u)=\tfrac12e^{-cu}$, avec $c=1$ pour A1 et B1, $c=1/2$ pour A2,
$c=4/5$ pour A3. La troncature est $\I=[0{,}05,\,0{,}95]$ --- valeur que
l'on retrouve en reproduisant les écarts $0{,}186$, $0{,}107$, $0{,}158$
annoncés par le manuscrit.

À $b=0$ on a $h=1/2$, donc $M_h\xi_j(u)=\xi_j(\max(0,u-1/2))$ puisque
$\xi_j$ décroît, et le module \eqref{eq:omega} se calcule en fermé.

\begin{proposition}[Perte de gap à largeur de bande dégénérée]
\label{prop:b0}
\statut{A, recalculé indépendamment}
\[
 \Omega_{\mathrm{A1/B1}}(1/2)=0{,}137713026295563,\quad
 \Omega_{\mathrm{A2}}(1/2)=0{,}081176584085900,\quad
 \Omega_{\mathrm{A3}}(1/2)=0{,}117500459573997 .
\]
Rapportées aux écarts de population, ces valeurs consomment
respectivement $73{,}88\,\%$, $75{,}65\,\%$ et $74{,}58\,\%$ du gap.
Autrement dit, à $b=0$ les fenêtres de population ne conservent
qu'environ \textbf{un quart} du gap nominal.
\end{proposition}

\noindent
\emph{Vérification.} Ces trois valeurs ont été recalculées de façon
indépendante à partir des définitions du manuscrit, par quadrature
numérique~; l'accord porte sur les quinze chiffres significatifs.

\begin{apport}
Ceci transforme une observation empirique en explication. Le papier peut
remplacer «~la largeur de bande dégénérée n'est jamais compétitive~» par
une raison chiffrée. C'est aussi une justification \emph{a priori} du
choix de grille, ce qui est toujours préférable à une justification
\emph{a posteriori} par les résultats de simulation.
\end{apport}

\subsection{Les plages de largeur de bande qui survivent}

\begin{proposition}[Région admissible pour le criblage]
\label{prop:region-b}
\statut{A}
La condition spatiale et la faisabilité locale sont satisfaites dès que
$h_n\to0$ et $n\alpha_nh_n/\log(p_nn)\to\infty$, soit
\[
 0<b<1-a .
\]
En particulier la restriction $b<1/2$ du manuscrit est \emph{inutile}
pour le criblage. Toute cellule à $b>0$ de la grille affichée
($a\in\{.25,\dots,.50\}$, $b\in\{0,\dots,.40\}$, somme maximale
$0{,}90<1$) est couverte, y compris les neuf réglages du bloc
d'agrégation.
\end{proposition}

\noindent
Trois régions sont perdues~: $b=0$ (proposition~\ref{prop:b0})~;
$a+b\ge1$, où le compte d'extrêmes locaux cesse de diverger~; et, pour
des gaps rétrécissants $\Delta_{\min,n}\asymp n^{-d}$, il faut de surcroît
$r_n=o(\Delta_{\min,n})$, soit $d<b$.

\begin{remarque}[Il y a deux conditions $b<1/2$ différentes]
\label{rem:deux-b}
La restriction $b<1/2$ du manuscrit provenait d'un argument global de
type DKW, et elle est bien inutile pour le criblage. Mais une condition
$b<1/2$ \emph{d'origine différente} réapparaîtra au
chapitre~\ref{ch:inference}, issue du contrôle du reste de Hájek, et
celle-là est nécessaire pour l'inférence calibrée sur le maximum. Il faut
donc distinguer les deux dans le registre~: pour le criblage seul,
$0<b<1-a$~; pour l'inférence simultanée, en outre $b<1/2$.
\end{remarque}

\clearpage

\section{L'annexe des hypothèses primitives}\label{ch:primitif}

Le manuscrit possède une annexe qui dérive les conditions utilisées par
les théorèmes à partir d'hypothèses «~primitives~» portant sur $\gamma$,
les queues conditionnelles et les noyaux de fibre. La question posée à la
campagne était~: cette annexe peut-elle être \emph{réellement} primitive,
ou doit-elle porter à jamais des conditions énoncées sur les objets
projetés~? La réponse est nuancée, et instructive.

\subsection{Le quotient minimal : tout passe par une intégrale scalaire}

\begin{proposition}[Factorisation exacte de la queue projetée]
\label{prop:factorisation}
\statut{A}
Sous la représentation primitive au niveau logarithmique
\[
 \Pp(Y>e^t\mid U=x)=c(x)\,e^{-t/\gamma(x)}\{1+r(t,x)\},
\]
posons, pour une coordonnée $j$ et une valeur $u$,
$x=\iota_j(u,v)$, $\xi=\xi_j(u)$ et
$D(v)=\frac1{\gamma(x)}-\frac1\xi\ \ge0$. Alors
\begin{equation}
 \bar F_{j,u}(e^t)=e^{-t/\xi_j(u)}\,A_{j,u}(t),
 \qquad
 A_{j,u}(t)=\int e^{-tD(v)}\,c\{\iota_j(u,v)\}\{1+r(t,\iota_j(u,v))\}\,K_j(u,dv).
 \label{eq:facto}
\end{equation}
L'intégrale scalaire $A_{j,u}$ est le \textbf{quotient complet minimal}
de la famille primitive~: à recodage bijectif près, c'est par elle seule
que toute la famille de queues conditionnelles entre dans la queue
projetée.
\end{proposition}

\begin{idee}
Une fois l'indice projeté connu, la seule chose qui reste à savoir de la
famille conditionnelle est \emph{comment la masse s'accumule près du
déficit nul} --- c'est-à-dire près des configurations presque
maximisantes --- pondérée par l'échelle $c$ et corrigée par le reste $r$.
Rien d'autre ne survit à la projection. Ceci explique d'un coup pourquoi
$(\gamma,K)$ ne peuvent pas suffire~: ils ignorent $c$ et $r$.
\end{idee}

\subsection{Une impossibilité, et sa portée exacte}

\begin{theoreme}[Non-détermination par $(\gamma,K)$]
\label{thm:nondet}
\statut{I, A}
Il n'existe aucune caractérisation nécessaire et suffisante de la
variation régulière de second ordre projetée, ni de la régularité
horizontale à seuil fini, qui soit fonction des seuls $(\gamma,K)$ sur la
classe non restreinte des queues conditionnelles.
\end{theoreme}

\begin{proof}[Témoin]
Prenons $p=1$, la même loi de $U$, le même noyau trivial, et
$\gamma\equiv\xi>0$. Premier modèle~: $\bar F_1(e^t)=e^{-t/\xi}$ (Pareto
exact). Second modèle, après un raccord fini anodin~:
\[
 \bar F_2(e^t)=\exp\{-t/\xi+t^{-1/2}\sin t\}.
\]
Pour $t$ assez grand, la dérivée logarithmique de $\bar F_2$ est
négative, donc c'est bien une fonction de survie. Pour tout $a\in\R$
fixé,
\[
 \log\frac{\bar F_2(e^{t+a})}{\bar F_2(e^{t})}
 =-\frac a\xi+(t+a)^{-1/2}\sin(t+a)-t^{-1/2}\sin t
 \longrightarrow-\frac a\xi,
\]
donc le \emph{premier} ordre est identique. Mais les incréments valent
\[
 (t+a)^{-1/2}\sin(t+a)-t^{-1/2}\sin t
 =2t^{-1/2}\sin(a/2)\cos(t+a/2)+o(t^{-1/2}),
\]
qui \emph{oscillent}. Les incréments correspondant à $a=\pi$ et
$a=\pi/2$ n'ont aucun rapport asymptotique stable~; aucun auxiliaire
régulier éventuellement non nul ne peut donc convenir simultanément pour
tout $a$ réel. Les deux modèles ont même $\gamma$ et mêmes noyaux, et
appartiennent à des classes de second ordre différentes.
\end{proof}

\begin{remarque}[Ce que ce théorème ne dit pas]
\label{rem:portee-nondet}
Il exclut une caractérisation à partir du \emph{descripteur réduit}
$(\gamma,K)$. Il n'exclut \emph{pas}~: que la variation régulière de
second ordre projetée existe~; qu'on l'impose directement sur les
quantiles projetés~; qu'on la déduise d'hypothèses primitives
\emph{plus riches} faisant intervenir $c(x)$, $r(t,x)$, leurs dérivées et
des couplages horizontaux~; qu'une condition suffisante utile existe~; ni
qu'une procédure statistique puisse la tester sur données complètes.
Confondre non-détermination et impossibilité serait une faute.
\end{remarque}

\subsection{L'enrichissement qui suffit, et il est vérifiable}

\begin{theoreme}[Transfert primitif]
\label{thm:transfert}
\statut{A$^-$}
Supposons, outre la variation régulière du premier ordre imposée en
couche séparée~:
\begin{align}
 &\text{(P2-SO)}\quad R_0(t)\to0,\qquad tR_1(t)\to0,
   \qquad\text{où } R_1(t)=\sup_x|\partial_tr(t,x)|,\notag\\
 &\text{(P3-RV)}\quad H_{j,u}(z)=a_{j,u}\,z^{\kappa_{j,u}}L_{j,u}(1/z)\{1+o(1)\},
   \qquad \kappa_{j,u}>0,\notag\\
 &\text{(P4)}\quad \text{couplage horizontal de } 1/\gamma,\ \log c,\ \log(1+r),\notag
\end{align}
où $H_{j,u}$ désigne la masse de déficit réciproque pondérée par $c$.
Alors le quantile projeté satisfait
\[
 Q_j(\cdot,u)\in2\mathrm{RV}_{\xi_j(u),\,0}
 \Bigl(-\frac{\xi_j(u)\kappa_{j,u}}{\log(\cdot)}\Bigr),
\]
uniformément sur les fibres utilisées par le score.
\end{theoreme}

\begin{idee}
La factorisation \eqref{eq:facto} montre que $A_{j,u}$ est une
transformée de Laplace de la masse de déficit. Un théorème taubérien de
Karamata transforme la variation régulière de cette masse \emph{en zéro}
en variation régulière de sa transformée \emph{à l'infini}. Il ne reste
qu'à inverser pour passer de la survie au quantile.
\end{idee}

\subsection{Pourquoi \texorpdfstring{$\rho=0$}{rho=0} et non \texorpdfstring{$-\kappa$}{-kappa} : trois variables à ne pas confondre}

C'est le point le plus facile à manquer, et il décide de toute la théorie
en aval.

\begin{proposition}[Le paramètre de second ordre vaut zéro]
\label{prop:rho0}
\statut{A}
Sous (P3-RV) avec $\kappa>0$, la survie projetée et le quantile supérieur
projeté ont pour paramètre de second ordre standard $\rho=0$, et non
$-\kappa$. Le nombre $-\kappa$ est l'indice de variation régulière du
facteur de Laplace $B(T)$ \emph{en la variable log-seuil} $T=\log y$,
c'est-à-dire $B\in\RV_{-\kappa}$~; ce n'est pas le paramètre de second
ordre relatif aux changements \emph{multiplicatifs} $y\mapsto\lambda y$.
La codimension effective $\kappa$ ne survit que dans la constante de
l'auxiliaire~:
\[
 A^S(y)\sim-\frac{\kappa}{\log y},
 \qquad
 A^Q(s)\sim-\frac{\xi\kappa}{\log s}.
\]
\end{proposition}

\noindent
Le tableau suivant est le meilleur garde-fou. Trois objets, trois
variables, trois comportements~:
\begin{center}
\begin{tabular}{@{}lll@{}}
\toprule
objet & variable & comportement dominant\\
\midrule
$H(\delta)$, masse de déficit & $\delta\downarrow0$ &
  $\delta^{\kappa}L(1/\delta)$ \quad --- \emph{polynomial}\\
$B(T)$, facteur de Laplace & $T\to\infty$ & $T^{-\kappa}L(T)$\\
$\bar F(y)$, survie projetée & $y\to\infty$ &
  $y^{-1/\xi}(\log y)^{-\kappa}L(\log y)$ \quad --- \emph{logarithmique}\\
\bottomrule
\end{tabular}
\end{center}

\noindent
Un comportement \emph{polynomial} dans la variable de déficit devient
donc un comportement \emph{logarithmique} dans la variable de réponse.
C'est cette translation, et rien d'autre, qui condamne le régime diffus à
$\rho=0$ --- ce qui aura les conséquences du chapitre~\ref{ch:biais}.

\subsection{Vérification explicite pour A1}

\begin{proposition}[Exposants de coin pour A1]
\label{prop:kappa}
\statut{A, recalculé indépendamment}
Pour A1, avec $\gamma(u)=\tfrac12\exp(-\sum_{a\le4}u_a)$,
$U_a=\Phi(Z_a)$ et bloc actif gaussien AR(1) de corrélation $1/4$~:
en conditionnant sur $U_j=u$ et en posant
$S=\sum_{a\le4,a\ne j}U_a$, le déficit réciproque exact est
\[
 D=2e^u(e^S-1)=2e^u\!\!\sum_{a\ne j}U_a+O(\|U_{-j}\|_1^2),
\]
donc le théorème de coin gaussien s'applique en dimension libre $3$, avec
$\kappa_j=\bm1^\top\Omega_j\bm1$ où $\Omega_j$ est l'inverse de la
covariance \textbf{conditionnelle} de $Z_{-j}$ sachant $Z_j$. Il vient
\[
 \kappa_j=\frac{34}{15}\ \ (j=1,4),
 \qquad
 \kappa_j=\frac{41}{15}\ \ (j=2,3),
\]
et les survies projetées
\begin{align*}
 \bar F_{1,u}(y)&\asymp y^{-1/\xi_1(u)}(\log y)^{-34/15}
   (\log\log y)^{-11/30}\exp\{-\tfrac4{15}z_u\sqrt{2\log\log y}\},\\
 \bar F_{2,u}(y)&\asymp y^{-1/\xi_2(u)}(\log y)^{-41/15}
   (\log\log y)^{-2/15}\exp\{-\tfrac8{15}z_u\sqrt{2\log\log y}\}.
\end{align*}
\end{proposition}

\noindent
\emph{Vérification.} Les deux exposants ont été recalculés
indépendamment~: $34/15=2{,}266667$ et $41/15=2{,}733333$, avec accord
exact. Le point délicat est bien le conditionnement~: avec la covariance
\emph{marginale} on obtiendrait $11/5=2{,}2$, ce qui serait faux.

\begin{apport}
La classe (P2-SO)/(P3-RV)/(P4) \emph{contient les quatre modèles de
simulation}, avec vérification complète~: le théorème de coin gaussien
fournit (P3-RV), les dérivées du reste satisfont (P2-SO), les couplages
existants donnent (P4). L'annexe peut donc \emph{démontrer} au lieu de
\emph{supposer}. La division éditoriale propre est~: théorème principal
sous les conditions projetées directes les plus faibles~; annexe
primitive comme théorème de transfert rigoureux pour une sous-classe
naturelle large. Il ne faut pas la présenter comme une équivalence
universelle --- elle ne l'est pas (remarque~\ref{rem:portee-nondet}).
\end{apport}

\subsection{La branche atomique : trois cas, pas deux}

\begin{proposition}[Classification des coins]
\label{prop:atome}
\statut{A}
\begin{center}
\begin{tabular}{@{}lll@{}}
\toprule
cas & $H$ près de zéro & correction projetée dominante\\
\midrule
coin continu & $H(0)=0$, $H(z)\sim az^\kappa L(1/z)$ &
  $A^Q(s)\sim-\xi\kappa/\log s$, \ $\rho=0$\\
atome $+$ coin continu & $H(z)=h_0+az^\beta L(1/z)$ &
  $A^Q(s)\asymp-(\log s)^{-\beta-1}L(\log s)$\\
atome $+$ trou spectral & $H(z)=h_0$ sur $(0,\delta_0)$ &
  $A^Q(s)\asymp s^{-\xi\delta_0}$, \ $\boxed{\rho=-\xi\delta_0<0}$\\
\bottomrule
\end{tabular}
\end{center}
Un mélange d'atome et de coin continu n'est \emph{pas} couvert par la
condition sans atome, mais l'est en appliquant la variation régulière à
$H(z)-H(0)$. Un mélange atomique \emph{non structuré}, pour lequel aucune
forme n'est imposée à $H(z)-H(0)$, n'a aucune correction de second ordre
identifiée et relève du seul premier ordre.
\end{proposition}

\begin{apport}
L'ancienne condition de doublement du manuscrit autorisait les atomes,
mais n'identifiait pas le taux résiduel, donc ne pouvait pas identifier
l'auxiliaire. La nouvelle formulation répare exactement cela. La
troisième ligne est à retenir~: c'est le seul endroit de toute la
campagne qui produise un $\rho$ \emph{strictement négatif}, et donc le
seul où la théorie classique de second ordre s'applique --- voir
section~\ref{sec:trou}.
\end{apport}

\clearpage

\section{La barrière du biais}\label{ch:biais}

Ce chapitre contient le principal résultat négatif de la campagne. Il
explique pourquoi tous les théorèmes du papier sont, et resteront,
confinés aux écarts \emph{fixes}.

\subsection{Pourquoi \texorpdfstring{$\rho=0$}{rho=0} interdit tout certificat polynomial}

\begin{theoreme}[Impossibilité d'un certificat de biais polynomial]
\label{thm:bias-imp}
\statut{I, A}
Sous la classe de second ordre projetée du manuscrit, il n'existe aucun
certificat de biais de queue d'ordre $o(n^{-d})$ pour $d>0$. Aucun
estimateur ne peut atteindre un taux polynomial uniforme sur cette
classe.
\end{theoreme}

\begin{idee}
Dans le régime de projection diffuse, $\rho=0$ et l'auxiliaire est
d'ordre $1/\log t$. Or la variation régulière de second ordre à $\rho=0$
affirme seulement que le reste est $o\{A(t)\}$~: elle \textbf{ne fournit
aucun taux} pour ce petit-$o$. On peut donc construire, à l'intérieur de
la même classe, un reste qui décroît plus lentement que \emph{toute}
puissance $t^{-\beta}$. Les corrections logarithmiques finies ne réparent
rien, puisque le défaut n'est pas dans le terme dominant mais dans
l'absence de contrôle du reste.
\end{idee}

\noindent
Le manuscrit énonce d'ailleurs déjà, sans en tirer la conséquence, une
condition (E2) qui donne un biais d'ordre $1/\log(1/\alpha)$ et ne couvre
que les écarts vérifiant $\Delta_{\min}\log(1/\alpha)\to\infty$~; avec
$\alpha=n^{-a}$, cela s'écrit $\Delta_{\min}\log n\to\infty$. Aucun
$\Delta_{\min}\asymp n^{-d}$ ne satisfait cela.

\begin{apport}
C'est une limite \emph{de la classe}, pas un défaut de preuve. L'article
doit donc énoncer explicitement que ses théorèmes valent à écarts fixes,
et \emph{pourquoi} --- au lieu de laisser croire que l'extension aux
écarts rétrécissants serait un exercice de routine. C'est un
renforcement de la rédaction, pas un aveu de faiblesse~: on remplace un
silence par un théorème.
\end{apport}

\subsection{La classe renforcée, et la bonne construction}

Un certificat polynomial redevient possible sous une hypothèse
\emph{quantitative de troisième ordre}, que l'on notera $\mathrm{QLH}_\beta$~:
tout le biais de projection non polynomial doit appartenir à une forme de
nuisance unidimensionnelle connue, avec un reste polynomial.

Sous $\mathrm{QLH}_\beta$, la construction statistiquement préférable
n'est \emph{pas} la séparation en puissance des seuils.

\begin{theoreme}[Extrapolation à rapport fixe]
\label{thm:extrapolation}
\statut{V}
Sous $\mathrm{QLH}_\beta$, l'estimateur combinant deux fractions de queue
différant d'un \emph{facteur constant} $\lambda$, avec des poids d'ordre
$\log(1/\alpha)$, atteint
\[
 \text{biais de queue}=O\{\alpha^\beta\log(1/\alpha)\},
 \qquad
 \text{inflation de variance}=O\{\log^2(1/\alpha)\}.
\]
Précisément, pour une fenêtre Pareto exacte gelée avec
$k_0/k_1\to\lambda$, la représentation de Rényi donne
$\Var(H_{m,k_0})=\xi^2/k_0$, $\Var(H_{m,k_1})=\xi^2/k_1$ et
$\Cov=\xi^2/k_0$, d'où
\[
 \Var(H^{\mathrm{BR}}_m)=\frac{\xi^2}{k_0}\{1+(\lambda-1)w_1^2\},
 \qquad
 v_{\lambda,\alpha}\sim\frac{\lambda-1}{(\log\lambda)^2}\log^2(1/\alpha).
\]
\end{theoreme}

\begin{proposition}[Rapport optimal]
\label{prop:lambda}
\statut{A, recalculé indépendamment}
La fonction $\lambda\mapsto(\lambda-1)/(\log\lambda)^2$ est minimisée en
\[
 \log\lambda^\star=1{,}59362,\qquad
 \lambda^\star=4{,}92155,\qquad
 \text{valeur minimale}=1{,}54414 .
\]
Le choix de $\lambda$ change les constantes, jamais la région
d'exposants.
\end{proposition}

\noindent
\emph{Vérification.} Ces trois valeurs ont été recalculées
indépendamment, avec accord exact sur les six chiffres.

\begin{idee}
L'inflation logarithmique est le prix, incompressible, d'identifier
\emph{faiblement} le coefficient d'un biais en $1/\log t$ à partir de
deux seuils voisins. Deux seuils trop proches rendent le système mal
conditionné (d'où $(\log\lambda)^{-2}$)~; deux seuils trop éloignés
placent le second trop loin dans la queue (d'où $\lambda-1$). L'optimum
est le compromis, et il vaut environ $5$.
\end{idee}

\subsection{La région admissible, et elle est non vide}

\begin{theoreme}[Région $(a,b,d)$ sous $\mathrm{QLH}_\beta$]
\label{thm:region}
\statut{V}
Avec $\alpha_n=n^{-a}$, $h_n=n^{-b}/2$, $p_n\le n^C$ et
$\Delta_{\min,n}\asymp n^{-d}$, les quatre contraintes sont
\begin{center}
\begin{tabular}{@{}ll@{}}
\toprule
contrainte & origine\\
\midrule
$a+b<1$ & faisabilité locale\\
$2d<1-a$ & erreur stochastique\\
$d<a\beta$ & certificat de biais de queue\\
$d<b$ & biais spatial et horizontal\\
\bottomrule
\end{tabular}
\end{center}
Un $a$ admissible existe si et seulement si $d/\beta<a<1-2d$, c'est-à-dire
si et seulement si
\[
 \boxed{\ d<\frac{\beta}{2\beta+1}\ }
\]
et l'on peut alors choisir $d<b<1-a$. La région est donc non vide~: pour
$\beta=1$, tout $d<1/3$ est atteignable.
\end{theoreme}

\subsection{Le prix : les quatre modèles sont exclus}

\begin{proposition}[A1--A3 et B1 ne satisfont pas $\mathrm{QLH}_\beta$]
\label{prop:modeles-exclus}
\statut{V}
La masse de déficit en coin gaussien des quatre modèles engendre, après
transformation de Laplace et inversion de quantile, des facteurs en
$\sqrt{\log\log t}$, $\log\log t/\log t$ et puissances inverses de
$\log\log t$. La correction par séparation en puissance y laisse des
biais d'ordre
\[
 \frac1{\log t\,(\log\log t)^{3/2}}
 \quad\text{ou}\quad
 \frac1{\log t\,(\log\log t)^{2}},
\]
à coefficients génériquement non nuls. Comme
$n^{d}/\{\log n(\log\log n)^\nu\}\to\infty$ pour tout $d>0$, ces résidus
ne sont \emph{jamais} $o(n^{-d})$. Les quatre modèles ne supportent donc
aucun théorème à écarts polynomialement rétrécissants.
\end{proposition}

\begin{remarque}[Compatibilité avec le chapitre~\ref{ch:primitif}]
Il n'y a pas de contradiction avec la
proposition~\ref{prop:kappa}, qui établit que ces mêmes modèles
satisfont (P3-RV) avec $\kappa>0$. Les deux énoncés portent sur des
variables différentes, comme le rappelle le tableau de la
proposition~\ref{prop:rho0}~: le facteur à variation lente absorbe toute
la structure en $\log\log$, qui est \emph{subordonnée} à l'auxiliaire
dominante en $1/\log s$. Les deux résultats tiennent simultanément.
\end{remarque}

\begin{apport}
Deux options honnêtes s'offrent à l'article. \emph{(i)} Énoncer le
théorème à écarts rétrécissants sous $\mathrm{QLH}_\beta$ et dire
franchement qu'aucun des quatre modèles ne l'illustre. \emph{(ii)}
Ajouter un \textbf{cinquième modèle de simulation}, construit pour
satisfaire $\mathrm{QLH}_\beta$ avec $\beta$ explicite, afin que le
théorème ait un témoin numérique. La seconde est nettement préférable, et
c'est un travail borné.
\end{apport}

\subsection{Une porte de sortie possible : le trou spectral}\label{sec:trou}

\begin{remarque}[Piste ouverte]
L'impossibilité du théorème~\ref{thm:bias-imp} porte sur la branche
$\rho=0$. Or la troisième ligne de la
proposition~\ref{prop:atome} --- atome au maximum avec trou spectral ---
donne $\rho=-\xi\delta_0$ \emph{strictement négatif}, avec une auxiliaire
qui est une vraie puissance. La théorie classique de second ordre s'y
applique, et le certificat polynomial devrait donc y être disponible sans
hypothèse de troisième ordre. Reste à déterminer la taille minimale
d'atome requise, et surtout si un trou $\delta_0\to0$ avec $n$
\emph{interpole} entre les deux régimes~: cela fournirait la frontière
exacte entre biais logarithmique et biais polynomial, c'est-à-dire la
condition structurelle précise sous laquelle le criblage à écarts
rétrécissants est possible. Cette question est ouverte à la date de ce
dossier.
\end{remarque}

\clearpage

\section{L'inférence simultanée calibrée sur le maximum}\label{ch:inference}

C'est le résultat central de la campagne, et celui qui a demandé le plus
d'étapes~: la question a été ouverte très tôt, déclarée close par la
négative, puis rouverte, puis résolue positivement.

\subsection{Ce que l'on veut, et pourquoi ce n'est pas gratuit}

L'article, en l'état, contrôle la multiplicité par Benjamini--Yekutieli,
valide sous dépendance arbitraire mais conservateur d'un facteur
$\log p$. Ce que l'on voudrait~: une valeur critique \emph{calibrée sur
le maximum}, obtenue par bootstrap à multiplicateurs gaussiens, avec la
vraie dépendance entre coordonnées.

Le passage exige une \emph{représentation linéaire par ligne}~: écrire le
score de chaque coordonnée comme une moyenne sur les lignes d'une
fonction d'influence de cette ligne seule, plus un reste négligeable
uniformément sur $p$ après multiplication par le taux d'approximation
gaussienne. Les lignes étant alors indépendantes, un multiplicateur
gaussien à \emph{ligne commune} rééchantillonne les $p$ scores
simultanément, et une borne de Chernozhukov--Chetverikov--Kato sur les
hyper-rectangles fournit la valeur critique.

\begin{idee}
La difficulté annoncée était que le score est bâti sur des \emph{rangs
empiriques} --- une ligne déplace le rang de toutes les autres --- et sur
une fenêtre dont l'appartenance est elle-même déterminée par les rangs.
L'influence d'une ligne n'est donc pas locale. Il s'est avéré que ce
n'était pas le vrai obstacle.
\end{idee}

\subsection{La fonction d'influence existe, et les rangs y contribuent deux termes de bord}

\begin{proposition}[Influence de ligne du fonctionnel de rang]
\label{prop:influence}
\statut{A}
Le fonctionnel de rang de population admet l'influence de ligne explicite
\begin{multline}
 \varphi_{j,n}(z)=\frac1{|\I|}\int_{\I}\frac1{2h_n}
 \Bigl[\1\{u-h_n<v_j\le u+h_n\}\,\chi_{j,u}(x)\\
 +\mu_{j,u}(u+h_n)\{u+h_n-\1(v_j\le u+h_n)\}
 -\mu_{j,u}(u-h_n)\{u-h_n-\1(v_j\le u-h_n)\}\Bigr]du .
 \label{eq:influence}
\end{multline}
Le premier terme est l'influence de Hill ordinaire~; les deux derniers
sont la \emph{correction de rang linéaire}, engendrée par le déplacement
des quantiles marginaux empiriques.
\end{proposition}

\noindent
Autrement dit~: la perturbation non locale des rangs se linéarise
parfaitement, et sa contribution est un couple de termes de bord. Ce
n'était pas l'obstacle.

\subsection{Une fausse impossibilité, et pourquoi elle était fausse}

Un premier verdict concluait à l'impossibilité, au motif qu'aucune
représentation ne peut avoir une influence \emph{bornée par une constante
$O(1)$}~: dans le modèle nul Pareto exact, la variance de l'influence est
d'ordre $1/\alpha_n$, et une influence bornée contredirait la borne
inférieure de fluctuation $1/\sqrt{n\alpha_n}$ déjà démontrée.

Ce raisonnement est valide, mais il vise une exigence dont personne n'a
besoin. Les bornes de
Chernozhukov--Chetverikov--Kato ne requièrent pas de sommants bornés~:
elles valent pour des enveloppes non bornées sous condition de norme
sous-exponentielle, au prix d'un taux explicite.

\begin{proposition}[L'étape gaussienne a de la marge]
\label{prop:cck}
\statut{A}
Après standardisation, l'influence de Hill a une enveloppe
sous-exponentielle $B_n\lesssim\alpha_n^{-1/2}$. Avec $\alpha_n=n^{-a}$,
$h_n=n^{-b}$, $p_n\le n^C$~:
\[
 \delta^{(6)}_{\mathrm{CCK},n}
 =O\Bigl[\Bigl\{\frac{\log^7n}{n^{1-a}}\Bigr\}^{1/6}\Bigr],
 \qquad
 \delta^{(4)}_{\mathrm{CCK},n}
 =O\Bigl[\Bigl\{\frac{\log^5(p_nn)}{n\alpha_n}\Bigr\}^{1/4}\Bigr],
\]
tous deux $o(1)$ si et seulement si $a<1$. La largeur de bande
\textbf{n'entre pas} dans cette approximation à sommants linéaires~; elle
ne revient que dans la faisabilité locale $n\alpha_nh_n\gg\log(p_nn)$,
soit $a+b<1$.

Pour une dimension véritablement sous-exponentielle,
$\log p_n\asymp n^\kappa$, les bornes informatives sont
$\kappa<(1-a)/7$ et, avec le théorème amélioré, $\kappa<(1-a)/5$.
\end{proposition}

\noindent
Dans toute la région de criblage du manuscrit, l'étape gaussienne n'est
donc pas l'étape limitante.

\subsection{Le vrai obstacle, et un faux pas instructif}

Restait un seul énoncé à prouver~: la négligeabilité de la somme de
toutes les composantes de Hoeffding d'ordre $\ge2$ --- l'erreur de
projection de Hájek. Le télescopage d'édition de rangs, déjà établi pour
la robustesse, contrôle une différence de \emph{premier} remplacement, et
ne dit rien des composantes dégénérées.

\begin{remarque}[L'énoncé faux qu'il faut connaître]
\label{rem:faux-pas}
Une intuition naturelle --- et fausse --- consiste à dire~: toute
différence seconde mixte gagne un facteur $(nh)^{-1}$ sur une différence
première, puisque remplacer une ligne ne déplace le rang de chaque autre
ligne que d'\emph{au plus une position}, de sorte que le second
remplacement voit une fenêtre quasi inchangée.

Cet énoncé est faux pour Hill, parce que \textbf{deux lignes peuvent
déplacer conjointement le seuil local}. L'intuition est
«~directionnellement juste mais fausse trajectoire par trajectoire~».
Aucune borne \emph{déterministe} de différence mixte ne peut fonctionner
ici, et il est utile de le savoir avant de perdre du temps à en chercher
une.
\end{remarque}

\subsection{La preuve : cinq étapes, toutes probabilistes}

\begin{theoreme}[Reste dégénéré uniformément négligeable]
\label{thm:hajek}
\statut{V}
Pour le score glissant local de Hill sur rangs empiriques, la somme de
toutes les composantes de Hoeffding d'ordre $\ge2$ est négligeable,
uniformément sur les $p_n$ coordonnées, à l'échelle exigée par
l'approximation gaussienne, sur la région
\[
 \boxed{\ 0<a<1,\qquad 0<b<\tfrac12,\qquad a+b<1 .\ }
\]
La représentation exacte par ligne a pour influence
\[
 \varphi_{j,n}(Z_i)=n\,b_{j,n}(U_{ij})\,
   h_{m,k}\bigl[-\log\{1-F_{Y|j}(Y_i\mid U_{ij})\}\bigr].
\]
\end{theoreme}

\begin{idee}
Puisque aucune borne déterministe ne marche, il faut une preuve
probabiliste. On linéarise chaque statistique de Hill locale à son
quantile intermédiaire \emph{déterministe} --- et non au quantile
empirique, ce qui serait circulaire --- puis on contrôle séparément deux
restes~: celui de la linéarisation locale, et celui de l'assignation des
rangs. Le coloriage du chapitre~\ref{ch:h} resert pour moyenner les
restes locaux.
\end{idee}

\noindent
\emph{Architecture de la preuve.}
\begin{enumerate}[leftmargin=2em,itemsep=3pt]
\item \textbf{Linéarisation conditionnelle} de chaque Hill local à son
quantile intermédiaire déterministe.
\item \textbf{Reste de Hájek local}~: on montre qu'il a une échelle de
queue $k^{-3/4}$, où $k\asymp n\alpha h$. Cela se décompose en trois
lemmes~: un reste de Bahadur local pour Hill~; la correction de
projection, d'ordre $O(k^{-1})$~; et une queue \emph{exponentielle} pour
le reste dégénéré local exact.
\item \textbf{Moyennage} de ces restes locaux au moyen du coloriage de
fenêtres disjointes.
\item \textbf{Statistique de rang linéaire} restante~: linéarisation, et
borne exponentielle sur son reste d'assignation de rang.
\item \textbf{Combinaison} des deux bornes exponentielles, puis union sur
$p_n$.
\end{enumerate}

\noindent
\emph{Sur la région d'exposants.} Les cinq parties polynomiales de la
borne sont respectivement d'ordre $n^{1-a-b}$, $n^{1-a+b}$, $n^{1-a}$,
$n^{1-2b}$, $n^{2-a-2b}$~; toutes sont contrôlées dès que $a+b<1$ et
$b<1/2$. Le repère $a+2b<1$ envisagé initialement est strictement plus
fort, donc suffisant.

\begin{remarque}[Les neuf réglages du manuscrit sont couverts]
Pour $a\in\{0{,}30,\,0{,}35,\,0{,}40\}$ et
$b\in\{0{,}10,\,0{,}15,\,0{,}20\}$, on a $\max(a+2b)=0{,}80<1$ et
$\max b=0{,}20<1/2$. Ce sont exactement les neuf réglages du bloc
d'agrégation. La théorie couvre donc la procédure telle qu'elle est
employée dans les simulations, ce qui est rarement le cas.
\end{remarque}

\noindent
Voir aussi la remarque~\ref{rem:deux-b}~: le $b<1/2$ qui apparaît ici est
d'origine différente de celui qui a été éliminé au
chapitre~\ref{ch:spatial}, et celui-ci est nécessaire.

\subsection{La conséquence}

\begin{corollaire}[Calibration sur le maximum, sans changer d'estimateur]
\label{cor:calibration}
\statut{V}
Soit $\psi_{j,n}=\varphi_{0,n}-\varphi_{j,n}$ l'\emph{influence d'écart},
différence entre l'influence de la ligne sur la ligne de base
inconditionnelle et son influence sur la coordonnée $j$. Alors,
uniformément en $j$,
\[
 \widehat\Delta_j-\Delta_{j,n}
 =\frac1n\sum_{i=1}^n\psi_{j,n}(Z_i)
 +o_\Pp\Bigl(\frac1{\sqrt{n\alpha\log p}}\Bigr),
\]
et la statistique à multiplicateur gaussien de \emph{ligne commune}
\[
 T^\ast_n=\max_{j\le p}\Bigl|\frac{\sqrt{\alpha}}{\widehat\sigma_j\sqrt n}
 \sum_{i=1}^ne_i\{\widehat\psi_{ij,n}-\overline{\widehat\psi}_{j,n}\}\Bigr|,
\]
où un seul multiplicateur $e_i$ est partagé par toutes les coordonnées,
a pour erreur d'approximation
$O[\{\log^5(pn)/(n\alpha)\}^{1/4}]$ plus l'erreur d'estimation de
covariance.
\end{corollaire}

\noindent
Le partage de $e_i$ est essentiel~: c'est lui qui préserve intégralement
la dépendance entre coordonnées, laquelle provient du fait que les mêmes
lignes, et les mêmes réponses extrêmes, alimentent plusieurs coordonnées.

\begin{apport}
C'est le gain le plus important pour l'article, et il est de nature
qualitative, non quantitative~: il fait passer d'un contrôle
\emph{conservateur} de la multiplicité à une \emph{calibration exacte}.
Trois points méritent d'être soulignés dans la rédaction.

\smallskip
\emph{Premièrement}, aucune scission d'échantillon n'est nécessaire. Une
construction alternative par échantillon scindé avait été proposée~; elle
a été retirée, car elle change l'estimateur, perd une fraction constante
de l'échantillon, et n'est pas démontrablement préférable au score brut.

\smallskip
\emph{Deuxièmement}, l'estimateur ne change pas. Le score calibré est
celui qui est déjà implémenté et déjà simulé.

\smallskip
\emph{Troisièmement}, la région couverte contient les réglages
effectivement utilisés, ce qui permet d'écrire une phrase que peu
d'articles peuvent écrire~: la théorie s'applique à la procédure telle
qu'elle a été exécutée.
\end{apport}

\begin{remarque}[Ce qui reste à vérifier]
Le théorème~\ref{thm:hajek} est marqué \textbf{V}~: l'audit adverse est
lancé et porte sur six points --- l'origine de l'exposant $3/4$~; la
légitimité de réutiliser le coloriage pour des restes \emph{dégénérés}
(dans une classe de couleur, les fenêtres sont disjointes en positions de
rang, mais les lignes sous-jacentes sont partagées entre coordonnées, de
sorte que l'indépendance pourrait n'être qu'intra-coordonnée)~; l'erreur
de transfert du quantile déterministe au quantile empirique~; la borne au
bord de troncature, où une fenêtre est rognée~; l'origine du $b<1/2$~; et
le caractère atteignable des conditions de cohérence de covariance.
Aucun de ces points ne doit être considéré comme acquis avant verdict.
\end{remarque}

\clearpage

\section{Robustesse, contamination, bornes inférieures}\label{ch:robust}

Ce chapitre rassemble les résultats acquis lors des premières vagues, dont
les preuves complètes figurent dans les documents annexes de la campagne.
On en donne ici l'énoncé, l'idée, et l'apport.

\subsection{Le télescopage d'édition de rangs}

\begin{theoreme}[Robustesse sous remplacement de lignes complètes]
\label{thm:telescopage}
\statut{A}
Sous remplacement adaptatif de lignes complètes, le score glissant borné
sur rangs reste robuste.
\end{theoreme}

\begin{idee}
La preuve initiale affirmait qu'une ligne modifiée n'altère que $O(nh)$
fenêtres. C'est faux~: une ligne modifiée déplace le rang de \emph{toutes}
les autres, donc affecte potentiellement toutes les fenêtres. La preuve
correcte est un \emph{télescopage}~: on remplace les lignes contaminées
une à une, et l'on borne chaque différence successive. Chaque
remplacement individuel ne déplace chaque rang que d'une position, donc
son effet sur le score est petit~; la somme des $\varepsilon_cn$
différences donne la borne. Le résultat survit, la preuve devait être
remplacée.
\end{idee}

\begin{apport}
Le \emph{résultat} du manuscrit est conservé, mais sa \emph{preuve} doit
être remplacée, avec deux ajustements~: plafond de winsorisation
\textbf{fixe} (un plafond croissant n'a pas de constante de contamination
indépendante de ce plafond), et départage déterministe des ex æquo. Il ne
faut pas écrire qu'aucune preuve n'est à jeter.
\end{apport}

\subsection{Le plancher de contamination est linéaire}

\begin{theoreme}[Plancher $\varepsilon_c/\alpha$]
\label{thm:contamination}
\statut{A}
Le plancher de contamination du score est d'ordre
$\varepsilon_c/\alpha$, \emph{linéaire} et non en racine.
\end{theoreme}

\begin{idee}
Deux voisinages de Huber d'intensité $\varepsilon$ se coupent avec
$\TV(P,Q)\le\varepsilon/(1-\varepsilon)$. Comme la distance en variation
totale exacte entre deux exponentielles de rapport $z$ vaut
$z(1-z)^{(1-z)/z}$, qui est du \emph{premier} ordre en $1-z$, le plancher
hérite de cette linéarité. Une contamination de fraction $\varepsilon_c$
concentrée dans la queue voit sa force amplifiée d'un facteur
$1/\alpha$, puisque seule la fraction $\alpha$ des observations
contribue.
\end{idee}

\begin{apport}
Un plancher \emph{linéaire} est bien plus contraignant qu'un plancher en
racine~: à $\alpha=n^{-0{,}3}$ et $n=2000$, le facteur d'amplification
est d'environ $10$. C'est une correction importante, car une version
antérieure de ce dossier annonçait $\sqrt{\varepsilon_c/\alpha}$, ce qui
était trop optimiste.
\end{apport}

\subsection{Bornes inférieures}

\begin{theoreme}[Transition en grande dimension par Fano avec liste]
\label{thm:fano}
\statut{A}
Sur un sous-modèle de queue plongé, la frontière de détectabilité est
\[
 n\alpha\,\Delta^2\asymp\log(p/d),
\]
avec un sélecteur de type top-$d$. L'appariement avec la borne supérieure
est \textbf{en taux, non en constante}.
\end{theoreme}

\begin{remarque}
L'argmin ne suffisait que pour $d=1$~; il faut un sélecteur de liste. Et
l'appariement à la constante près n'a pas été établi --- ne pas l'écrire.
\end{remarque}

\begin{proposition}[Agrégation par rang minimal]
\label{prop:agregation}
\statut{A}
L'agrégation sur $M$ réglages voisins n'est \emph{ni} un oracle \emph{ni}
une intersection~: c'est une couverture par \emph{union des tops}, avec
inflation de liste. Si un seul réglage est bon, la garantie porte sur le
top $M\cdot s_D$.
\end{proposition}

\begin{apport}
Avec $M=9$ et $s_D=4$, la garantie porte sur le top $36$, et non sur le
top $20$. Le manuscrit annonce «~Sure-20~»~: il faut soit rapporter
Sure-36, soit énoncer explicitement que Sure-20 est une performance
empirique et non la quantité couverte par le théorème. C'est le genre de
décalage qu'un rapporteur relève systématiquement.
\end{apport}

\clearpage

\section{Ce qu'il faudrait changer dans l'article}\label{ch:changements}

Ce chapitre traduit le dossier en instructions de rédaction.

\subsection{Corrections obligatoires}

\begin{enumerate}[leftmargin=2em,itemsep=4pt]
\item \textbf{Abrégé et théorème de score}~: remplacer
$\sqrt{\log(pn)/(n\alpha h)}$ par $\sqrt{\log p/(n\alpha)}$, et la
condition de criblage $\log(pn)=o(n\alpha h\Delta^2)$ par
$\log p=o(n\alpha\Delta^2)$. Conserver $n\alpha h$ pour la faisabilité
locale, le biais spatial et la concentration ponctuelle, en le disant.
\item \textbf{Cible}~: remplacer $\max_{\supp}\gamma$ par
$\esssup\gamma$ partout~; viser $\D$ et non $A$~; présenter $\D=A$ comme
une hypothèse d'accessibilité explicite. L'hypothèse (S) telle
qu'imprimée n'identifie pas $A$.
\item \textbf{Condition spatiale}~: remplacer toute hypothèse de
continuité du profil par la condition quantitative
$\Omega_n(r_n)=o(\Delta_{\min})$, imposée au rayon gonflé $r_n$ et non à
$h_n$~; signaler que les fenêtres visent un \emph{supremum essentiel
mobile}, pas une moyenne.
\item \textbf{Portée des impossibilités}~: écrire «~non détectable
\emph{uniformément}~» pour l'apparition retardée~; restreindre
l'invisibilité marginale aux décisions mesurables par rapport au seul
couple $(U_j,Y)$~; ne pas présenter (P2) comme \emph{la} condition
nécessaire.
\item \textbf{Agrégation}~: top $M\cdot s_D=36$, pas $20$.
\item \textbf{Écarts fixes}~: énoncer explicitement que les théorèmes
valent à écarts fixes, et pourquoi (théorème~\ref{thm:bias-imp}).
\item \textbf{Exposants}~: écrire $0<a<1$ (le zéro était omis) et
$0<b<1-a$ pour le criblage~; ajouter $b<1/2$ pour l'inférence.
\item \textbf{Robustesse}~: remplacer la preuve par le télescopage~;
plafond de winsorisation fixe~; plancher $\varepsilon_c/\alpha$ linéaire.
\item \textbf{Application}~: statut de criblage exploratoire, non de
découverte certifiée~; hiérarchiser les prédicteurs en consensus /
stables / sensibles au réglage / promus par un seul réglage.
\item \textbf{Nombres}~: écrire systématiquement le design \emph{et} la
formule à côté de chaque valeur numérique. Les valeurs $0{,}67$ et
$0{,}611$ sont justes chacune pour son design (simulation et
application)~; leur ressemblance après correction est trompeuse.
\end{enumerate}

\subsection{Ajouts recommandés}

\begin{enumerate}[leftmargin=2em,itemsep=4pt]
\item \textbf{Section d'inférence simultanée} fondée sur le
chapitre~\ref{ch:inference}, remplaçant Benjamini--Yekutieli par une
valeur critique calibrée sur le maximum. C'est l'ajout le plus
substantiel.
\item \textbf{Annexe primitive réécrite} sous (P2-SO)/(P3-RV)/(P4),
démontrant au lieu de supposer, avec vérification pour les quatre
modèles.
\item \textbf{Explication de la colonne $b=0$} par la
proposition~\ref{prop:b0}, qui transforme une observation empirique en
raison chiffrée.
\item \textbf{Cinquième modèle de simulation} satisfaisant
$\mathrm{QLH}_\beta$, si l'on souhaite énoncer le théorème à écarts
rétrécissants avec un témoin numérique.
\item \textbf{Discussion des trois impossibilités}, qui justifient les
hypothèses au lieu de les subir.
\end{enumerate}

\section{Ce qui reste ouvert}

\begin{enumerate}[leftmargin=2em,itemsep=4pt]
\item \textbf{Audit du théorème~\ref{thm:hajek}} --- six points listés
en fin de chapitre~\ref{ch:inference}. C'est la priorité~: le résultat
central ne doit pas entrer dans l'article avant verdict.
\item \textbf{La branche à trou spectral} (section~\ref{sec:trou}) --- un
trou $\delta_0\to0$ interpole-t-il entre biais logarithmique et biais
polynomial~? Cela donnerait la frontière exacte du criblage à écarts
rétrécissants.
\item \textbf{Le régime intermédiaire} --- dans la classe primitive, le
biais dominant est d'ordre exactement $1/\log t$, à coefficient
\emph{identifié}. Peut-on l'estimer, le soustraire, et itérer jusqu'à
$1/\log^J$~? Cela donnerait un régime d'écarts strictement entre les
écarts fixes et les écarts polynomiaux.
\item \textbf{Rédaction complète} des résultats marqués \textbf{S}.
\item \textbf{Compression} --- réduire l'ensemble des preuves du papier
sous une forme sensiblement plus courte, sans affaiblir aucun énoncé.
\end{enumerate}

\end{document}
